\chapter{Budoucí vylepšení}

Současný prototyp plní základní funkci měření impedance, avšak několik oblastí nabízí prostor pro výrazné zlepšení. Následující podkapitoly navrhují možná budoucí vylepšení, která by zvýšila přesnost, použitelnost a uživatelský komfort zařízení.

\section{Kalibrace}

Přesnost měření je omezena systematickými chybami analogové cesty -- offsetem a ziskem A/D převodníků, fázovým posunem mezi kanály, tolerancemi referenčních bočníků a parazitními vlastnostmi měřicích vodičů. V současné podobě prototyp neimplementuje žádnou kalibrační proceduru, všechna měření jsou prováděna s hrubými hodnotami.

Implementace kalibrace by zahrnovala několik kroků:

\begin{itemize}
    \item \textbf{Kalibrace offsetu} -- měřením při nulovém vstupním signálu a následným odečtením střední hodnoty z měřených dat.
    \item \textbf{Kalibrace zisku} -- měřením referenčního rezistoru o známé hodnotě pro každý rozsah a výpočtem korekčního koeficientu.
    \item \textbf{Open a short kalibrace} -- kompenzací parazitní kapacity (open) a parazitního odporu a indukčnosti (short) měřicích vodičů a přípravku.
\end{itemize}

Tyto kalibrační kroky by ideálně probíhaly automaticky při spuštění přístroje nebo na vyžádání uživatelem. Kalibrační koeficienty by byly ukládány do paměti mikrokontroléru pro jednotlivé frekvenční i proudové rozsahy.

\section{Automatické přepínání rozsahů}

Při ručním nastavování rozsahu musí uživatel na PC aplikaci sledovat amplitudy signálů a podle nich volit vhodný bočník a zesílení. Tento proces je zdlouhavý a vyžaduje zkušenost.

Automatické přepínání rozsahů by fungovalo na principu předběžného měření: přístroj by nejprve provedl rychlé změření při nejcitlivějším rozsahu, vyhodnotil, zda nedochází k ořezání signálu nebo naopak k příliš nízké úrovni, a podle toho by zvolil optimální kombinaci bočníku a zesílení. Následně by provedl ostré měření s nastavenými parametry.

Tato funkce by vyžadovala rozšíření firmwaru o stavový automat a úpravu komunikačního protokolu. Vhodné by bylo také přidání hardwarového detektoru ořezání signálu, který by umožnil rychlejší reakci na přebuzení.

\section{Fyzické ovládací rozhraní}

Současný prototyp je ovládán výhradně přes PC aplikaci, což omezuje jeho použití v situacích, 
kde není k dispozici počítač. 
Přidání fyzického ovládacího rozhraní by výrazně zvýšilo samostatnost 
a praktickou použitelnost zařízení.

Prototyp obsahuje konektor pro rozšíření s vyvedenými sběrnicemi $I^{2}C$, SPI a UART.
Bylo zamýšleno, že se do něj připojí deska ovládacího panelu s vlastním mikrokontrolérem, 
mohl by se tak zachovat podobný komunikační protokol jako po sériové lince s PC.

\section{Elektrická ochrana}

Současný prototyp neobsahuje žádné obvody na ochranu proti elektrostatickému výboji, přepětí na vstupech ani galvanické oddělení. Pro nasazení v méně kontrolovaném prostředí by bylo vhodné tyto ochrany doplnit.

\subsection{Ochrana proti ESD}

Měřicí vodiče a BNC konektory jsou přímou cestou pro elektrostatický výboj do citlivé 
analogové části. Ochrana by mohla být realizována pomocí nízkokapacitních TVS diod 
 zapojených mezi signálové vodiče a zem. 
Tyto diody mají zanedbatelný vliv na měřený signál při běžném provozu, 
ale při ESD výboji odvedou proud do země dříve, 
než dojde k poškození operačních zesilovačů nebo ADC převodníku.

\subsection{Ochrana proti přepětí}

Připojení externího napětí na měřicí svorky (např. při měření na již napájeném obvodu) 
by mohlo zničit vstupní stupně zesilovačů. Vhodným doplňkem by byly zpětně zapojené diody 
k napájecím kolejnicím v kombinaci s PTC pojistkami omezujícími proud. 
Pro vyšší úroveň ochrany by bylo možné použít integrované přepěťové ochrany, 
které kombinují TVS diodu a proudové omezení v jednom pouzdře.

\subsection{Galvanické oddělení}

Protože je zařízení napájeno z USB portu počítače, vzniká zemní smyčka mezi přístrojem a 
počítačem. U měření velmi nízkých impedancí se tato smyčka může projevit 
jako přídavný šum nebo chyba měření. Galvanické oddělení sériové linky pomocí optoizolátoru 
(např. ISO7240) nebo integrovaného izolovaného rozhraní (např. ADuM1201) 
by tuto smyčku přerušilo. Zároveň by optoizolátor poskytl dodatečnou ochranu počítače 
v případě poruchy na měřeném zařízení.

